% Vietnamese translation for OI Public Library
% Source-File: IOI中国国家候选队论文/2003/王知昆/王知昆.ppt
\documentclass[11pt]{article}
\usepackage{oipl-vn}

\title{Bản trình chiếu: Hình chữ nhật con lớn nhất}
\author{Wang Zhikun}
\date{2003}

\begin{document}
\maketitle

\origtitle{Slide companion for maximum subrectangle problems}
\sourcepath{IOI China candidate team papers / 2003 / Wang Zhikun / Wang Zhikun.ppt}

\section{Phạm vi bản dịch}

Bản trình chiếu đi kèm bài viết về tư tưởng cực đại hóa trong bài toán hình
chữ nhật con. Các slide khôi phục được chứa tên John, các giới hạn \(5000\),
\(30000\), và các mảng \(H[i,j]\), \(L[i,j]\), \(R[i,j]\).

\section{Mô hình}

Với mỗi ô \((i,j)\), đặt \(H[i,j]\) là chiều cao liên tiếp có thể kéo lên từ
ô đó. Hai mảng \(L[i,j]\) và \(R[i,j]\) lưu biên trái, biên phải xa nhất mà
hình chữ nhật có chiều cao \(H[i,j]\) còn hợp lệ. Các công thức trong slide
có dạng:
\[
  H[i,j]=H[i-1,j]+1,
\]
\[
  L[i,j]=\max(L[i,j],L[i-1,j]),\qquad
  R[i,j]=\min(R[i,j],R[i-1,j]).
\]

\section{Độ phức tạp}

Các slide so sánh các cách làm \(O(S^5)\), \(O(S^3)\), \(O(S^2)\) và
\(O(S)\) cho từng bước phụ. Bằng cách duy trì chiều cao và biên hợp lệ theo
từng hàng, bài toán trên lưới \(N\times M\) có thể xử lý trong \(O(NM)\).

\section{Liên hệ với bài viết}

Bài viết đã dịch trình bày nhiều biến thể hơn. Bản trình chiếu tập trung vào
cách chuyển từ quét hình chữ nhật trực tiếp sang quy hoạch theo histogram và
các biên cực đại.

\end{document}
